% Options for packages loaded elsewhere
% Options for packages loaded elsewhere
\PassOptionsToPackage{unicode}{hyperref}
\PassOptionsToPackage{hyphens}{url}
\PassOptionsToPackage{dvipsnames,svgnames,x11names}{xcolor}
%
\documentclass[
  12pt,
  letterpaper,
]{article}
\usepackage{xcolor}
\usepackage[margin=1in]{geometry}
\usepackage{amsmath,amssymb}
\setcounter{secnumdepth}{-\maxdimen} % remove section numbering
\usepackage{iftex}
\ifPDFTeX
  \usepackage[T1]{fontenc}
  \usepackage[utf8]{inputenc}
  \usepackage{textcomp} % provide euro and other symbols
\else % if luatex or xetex
  \usepackage{unicode-math} % this also loads fontspec
  \defaultfontfeatures{Scale=MatchLowercase}
  \defaultfontfeatures[\rmfamily]{Ligatures=TeX,Scale=1}
\fi
\usepackage{lmodern}
\ifPDFTeX\else
  % xetex/luatex font selection
  \setmainfont[]{Times New Roman}
\fi
% Use upquote if available, for straight quotes in verbatim environments
\IfFileExists{upquote.sty}{\usepackage{upquote}}{}
\IfFileExists{microtype.sty}{% use microtype if available
  \usepackage[]{microtype}
  \UseMicrotypeSet[protrusion]{basicmath} % disable protrusion for tt fonts
}{}
\usepackage{setspace}
\makeatletter
\@ifundefined{KOMAClassName}{% if non-KOMA class
  \IfFileExists{parskip.sty}{%
    \usepackage{parskip}
  }{% else
    \setlength{\parindent}{0pt}
    \setlength{\parskip}{6pt plus 2pt minus 1pt}}
}{% if KOMA class
  \KOMAoptions{parskip=half}}
\makeatother
% Make \paragraph and \subparagraph free-standing
\makeatletter
\ifx\paragraph\undefined\else
  \let\oldparagraph\paragraph
  \renewcommand{\paragraph}{
    \@ifstar
      \xxxParagraphStar
      \xxxParagraphNoStar
  }
  \newcommand{\xxxParagraphStar}[1]{\oldparagraph*{#1}\mbox{}}
  \newcommand{\xxxParagraphNoStar}[1]{\oldparagraph{#1}\mbox{}}
\fi
\ifx\subparagraph\undefined\else
  \let\oldsubparagraph\subparagraph
  \renewcommand{\subparagraph}{
    \@ifstar
      \xxxSubParagraphStar
      \xxxSubParagraphNoStar
  }
  \newcommand{\xxxSubParagraphStar}[1]{\oldsubparagraph*{#1}\mbox{}}
  \newcommand{\xxxSubParagraphNoStar}[1]{\oldsubparagraph{#1}\mbox{}}
\fi
\makeatother

\usepackage{color}
\usepackage{fancyvrb}
\newcommand{\VerbBar}{|}
\newcommand{\VERB}{\Verb[commandchars=\\\{\}]}
\DefineVerbatimEnvironment{Highlighting}{Verbatim}{commandchars=\\\{\}}
% Add ',fontsize=\small' for more characters per line
\usepackage{framed}
\definecolor{shadecolor}{RGB}{241,243,245}
\newenvironment{Shaded}{\begin{snugshade}}{\end{snugshade}}
\newcommand{\AlertTok}[1]{\textcolor[rgb]{0.68,0.00,0.00}{#1}}
\newcommand{\AnnotationTok}[1]{\textcolor[rgb]{0.37,0.37,0.37}{#1}}
\newcommand{\AttributeTok}[1]{\textcolor[rgb]{0.40,0.45,0.13}{#1}}
\newcommand{\BaseNTok}[1]{\textcolor[rgb]{0.68,0.00,0.00}{#1}}
\newcommand{\BuiltInTok}[1]{\textcolor[rgb]{0.00,0.23,0.31}{#1}}
\newcommand{\CharTok}[1]{\textcolor[rgb]{0.13,0.47,0.30}{#1}}
\newcommand{\CommentTok}[1]{\textcolor[rgb]{0.37,0.37,0.37}{#1}}
\newcommand{\CommentVarTok}[1]{\textcolor[rgb]{0.37,0.37,0.37}{\textit{#1}}}
\newcommand{\ConstantTok}[1]{\textcolor[rgb]{0.56,0.35,0.01}{#1}}
\newcommand{\ControlFlowTok}[1]{\textcolor[rgb]{0.00,0.23,0.31}{\textbf{#1}}}
\newcommand{\DataTypeTok}[1]{\textcolor[rgb]{0.68,0.00,0.00}{#1}}
\newcommand{\DecValTok}[1]{\textcolor[rgb]{0.68,0.00,0.00}{#1}}
\newcommand{\DocumentationTok}[1]{\textcolor[rgb]{0.37,0.37,0.37}{\textit{#1}}}
\newcommand{\ErrorTok}[1]{\textcolor[rgb]{0.68,0.00,0.00}{#1}}
\newcommand{\ExtensionTok}[1]{\textcolor[rgb]{0.00,0.23,0.31}{#1}}
\newcommand{\FloatTok}[1]{\textcolor[rgb]{0.68,0.00,0.00}{#1}}
\newcommand{\FunctionTok}[1]{\textcolor[rgb]{0.28,0.35,0.67}{#1}}
\newcommand{\ImportTok}[1]{\textcolor[rgb]{0.00,0.46,0.62}{#1}}
\newcommand{\InformationTok}[1]{\textcolor[rgb]{0.37,0.37,0.37}{#1}}
\newcommand{\KeywordTok}[1]{\textcolor[rgb]{0.00,0.23,0.31}{\textbf{#1}}}
\newcommand{\NormalTok}[1]{\textcolor[rgb]{0.00,0.23,0.31}{#1}}
\newcommand{\OperatorTok}[1]{\textcolor[rgb]{0.37,0.37,0.37}{#1}}
\newcommand{\OtherTok}[1]{\textcolor[rgb]{0.00,0.23,0.31}{#1}}
\newcommand{\PreprocessorTok}[1]{\textcolor[rgb]{0.68,0.00,0.00}{#1}}
\newcommand{\RegionMarkerTok}[1]{\textcolor[rgb]{0.00,0.23,0.31}{#1}}
\newcommand{\SpecialCharTok}[1]{\textcolor[rgb]{0.37,0.37,0.37}{#1}}
\newcommand{\SpecialStringTok}[1]{\textcolor[rgb]{0.13,0.47,0.30}{#1}}
\newcommand{\StringTok}[1]{\textcolor[rgb]{0.13,0.47,0.30}{#1}}
\newcommand{\VariableTok}[1]{\textcolor[rgb]{0.07,0.07,0.07}{#1}}
\newcommand{\VerbatimStringTok}[1]{\textcolor[rgb]{0.13,0.47,0.30}{#1}}
\newcommand{\WarningTok}[1]{\textcolor[rgb]{0.37,0.37,0.37}{\textit{#1}}}

\usepackage{longtable,booktabs,array}
\usepackage{calc} % for calculating minipage widths
% Correct order of tables after \paragraph or \subparagraph
\usepackage{etoolbox}
\makeatletter
\patchcmd\longtable{\par}{\if@noskipsec\mbox{}\fi\par}{}{}
\makeatother
% Allow footnotes in longtable head/foot
\IfFileExists{footnotehyper.sty}{\usepackage{footnotehyper}}{\usepackage{footnote}}
\makesavenoteenv{longtable}
\usepackage{graphicx}
\makeatletter
\newsavebox\pandoc@box
\newcommand*\pandocbounded[1]{% scales image to fit in text height/width
  \sbox\pandoc@box{#1}%
  \Gscale@div\@tempa{\textheight}{\dimexpr\ht\pandoc@box+\dp\pandoc@box\relax}%
  \Gscale@div\@tempb{\linewidth}{\wd\pandoc@box}%
  \ifdim\@tempb\p@<\@tempa\p@\let\@tempa\@tempb\fi% select the smaller of both
  \ifdim\@tempa\p@<\p@\scalebox{\@tempa}{\usebox\pandoc@box}%
  \else\usebox{\pandoc@box}%
  \fi%
}
% Set default figure placement to htbp
\def\fps@figure{htbp}
\makeatother





\setlength{\emergencystretch}{3em} % prevent overfull lines

\providecommand{\tightlist}{%
  \setlength{\itemsep}{0pt}\setlength{\parskip}{0pt}}



 


\usepackage{booktabs}
\usepackage{longtable}
\usepackage{float}
\usepackage{setspace}
\doublespacing
% PAR heading styles
\usepackage{titlesec}
% Principal subheads: 14pt bold
\titleformat{\section}{\normalfont\fontsize{14}{16}\bfseries}{\thesection}{1em}{}
% Secondary subheads: 12pt bold
\titleformat{\subsection}{\normalfont\fontsize{12}{14}\bfseries}{\thesubsection}{1em}{}
% Sub-subheads: 12pt bold italic, run-in
\titleformat{\subsubsection}[runin]{\normalfont\fontsize{12}{14}\bfseries\itshape}{\thesubsubsection}{1em}{}[.]
% Abstract formatting
\renewenvironment{abstract}{\par\bigskip\noindent\textbf{Abstract}\par\smallskip\noindent}{\par\bigskip}
% Remove section numbers
\setcounter{secnumdepth}{0}
\makeatletter
\@ifpackageloaded{caption}{}{\usepackage{caption}}
\AtBeginDocument{%
\ifdefined\contentsname
  \renewcommand*\contentsname{Table of contents}
\else
  \newcommand\contentsname{Table of contents}
\fi
\ifdefined\listfigurename
  \renewcommand*\listfigurename{List of Figures}
\else
  \newcommand\listfigurename{List of Figures}
\fi
\ifdefined\listtablename
  \renewcommand*\listtablename{List of Tables}
\else
  \newcommand\listtablename{List of Tables}
\fi
\ifdefined\figurename
  \renewcommand*\figurename{Figure}
\else
  \newcommand\figurename{Figure}
\fi
\ifdefined\tablename
  \renewcommand*\tablename{Table}
\else
  \newcommand\tablename{Table}
\fi
}
\@ifpackageloaded{float}{}{\usepackage{float}}
\floatstyle{ruled}
\@ifundefined{c@chapter}{\newfloat{codelisting}{h}{lop}}{\newfloat{codelisting}{h}{lop}[chapter]}
\floatname{codelisting}{Listing}
\newcommand*\listoflistings{\listof{codelisting}{List of Listings}}
\makeatother
\makeatletter
\makeatother
\makeatletter
\@ifpackageloaded{caption}{}{\usepackage{caption}}
\@ifpackageloaded{subcaption}{}{\usepackage{subcaption}}
\makeatother
\usepackage{bookmark}
\IfFileExists{xurl.sty}{\usepackage{xurl}}{} % add URL line breaks if available
\urlstyle{same}
\hypersetup{
  pdftitle={Appendix B: Extended Methods},
  colorlinks=true,
  linkcolor={blue},
  filecolor={Maroon},
  citecolor={Blue},
  urlcolor={Blue},
  pdfcreator={LaTeX via pandoc}}


\title{Appendix B: Extended Methods}
\author{}
\date{}
\begin{document}
\maketitle


\setstretch{2}
\subsection{B.1 Cox Proportional Hazards
Model}\label{b.1-cox-proportional-hazards-model}

\subsubsection{Model Specification}\label{model-specification}

The primary model is the Cox proportional hazards regression:

\[
h(t | \mathbf{X}) = h_0(t) \exp(\mathbf{X}\boldsymbol{\beta})
\]

where:

\begin{itemize}
\tightlist
\item
  \(h(t | \mathbf{X})\) is the hazard of completion at time \(t\) given
  covariates \(\mathbf{X}\)
\item
  \(h_0(t)\) is the unspecified baseline hazard function
\item
  \(\boldsymbol{\beta}\) is the vector of regression coefficients
\end{itemize}

The semi-parametric nature of the Cox model allows estimation without
specifying the baseline hazard, reducing model misspecification risk.

\subsubsection{Proportional Hazards
Assumption}\label{proportional-hazards-assumption}

The Cox model assumes that hazard ratios are constant over time:

\[
\frac{h(t | X_1)}{h(t | X_2)} = \exp(\boldsymbol{\beta}(X_1 - X_2)) \quad \forall t
\]

We test this assumption using Schoenfeld residuals. For each covariate,
we regress scaled Schoenfeld residuals on time functions and test for
non-zero slopes. Non-significant tests (p \textgreater{} 0.05) support
the proportional hazards assumption.

\subsubsection{Partial Likelihood
Estimation}\label{partial-likelihood-estimation}

Parameters are estimated by maximizing the partial likelihood:

\[
L(\boldsymbol{\beta}) = \prod_{i: \delta_i = 1} \frac{\exp(\mathbf{X}_i \boldsymbol{\beta})}{\sum_{j \in R(t_i)} \exp(\mathbf{X}_j \boldsymbol{\beta})}
\]

where:

\begin{itemize}
\tightlist
\item
  \(\delta_i = 1\) indicates an observed event (completion)
\item
  \(R(t_i)\) is the risk set at time \(t_i\) (all subjects still at
  risk)
\end{itemize}

Standard errors are computed from the observed information matrix.

\subsubsection{Hazard Ratio
Interpretation}\label{hazard-ratio-interpretation}

Hazard ratios (HR) quantify the multiplicative change in completion
hazard per unit change in the covariate:

\begin{itemize}
\tightlist
\item
  HR \textgreater{} 1: Higher velocity associated with faster completion
\item
  HR \textless{} 1: Higher velocity associated with slower completion
\item
  HR = 1: No association
\end{itemize}

For velocity measured in percentage points per quarter, HR = 4.60 means
each additional percentage point of quarterly velocity is associated
with a 4.60-fold increase in the instantaneous completion hazard.

\subsection{B.2 Stratified Analysis}\label{b.2-stratified-analysis}

\subsubsection{Approach}\label{approach}

Context-specific effects are estimated by fitting separate Cox models
within subgroups:

\[
h_k(t | \mathbf{X}) = h_{0k}(t) \exp(\mathbf{X}\boldsymbol{\beta}_k) \quad \text{for } k = 1, \ldots, K
\]

where \(k\) indexes strata (disaster types, program types, experience
levels).

\subsubsection{Stratification Variables}\label{stratification-variables}

\begin{longtable}[]{@{}ll@{}}
\toprule\noalign{}
Variable & Categories \\
\midrule\noalign{}
\endhead
\bottomrule\noalign{}
\endlastfoot
Disaster Type & Hurricane, Wildfire, Flood, Other \\
Program Type & Housing, Infrastructure, Administration \\
Experience & Novice, Experienced \\
Program Phase & Early, Middle, Late \\
\end{longtable}

\subsubsection{Small Sample
Considerations}\label{small-sample-considerations}

Stratified analyses reduce sample sizes, potentially compromising
statistical power. We report confidence intervals alongside p-values to
convey uncertainty. Wide confidence intervals (e.g., HR=51.09, 95\% CI:
4.01-650.90 for wildfire) indicate effect magnitude is estimated
imprecisely even when the direction is statistically significant.

\subsection{B.3 Interaction Models}\label{b.3-interaction-models}

\subsubsection{Specification}\label{specification}

Formal tests of heterogeneity use interaction terms:

\[
\log h(t) = \log h_0(t) + \beta_1 \text{Velocity} + \beta_2 \text{Context} + \beta_3 (\text{Velocity} \times \text{Context})
\]

The interaction coefficient \(\beta_3\) tests whether velocity effects
differ significantly across contexts. Wald tests assess
\(H_0: \beta_3 = 0\).

\subsubsection{Power Considerations}\label{power-considerations}

Interaction tests require larger samples than main effect tests. With
N=152 and substantial censoring, power to detect interaction effects is
limited. Non-significant interactions should be interpreted cautiously;
they may reflect insufficient power rather than true homogeneity.

\subsection{B.4 Trajectory Clustering}\label{b.4-trajectory-clustering}

\subsubsection{K-Means Algorithm}\label{k-means-algorithm}

Velocity trajectories are clustered using K-means on standardized
quarterly velocity profiles:

\begin{enumerate}
\def\labelenumi{\arabic{enumi}.}
\tightlist
\item
  \textbf{Standardization}: Each grantee's velocity time series is
  z-scored to remove scale differences
\item
  \textbf{Alignment}: Series are interpolated to equal lengths (20 time
  points)
\item
  \textbf{Clustering}: K-means minimizes within-cluster sum of squared
  distances:
\end{enumerate}

\[
\min_{\mathbf{C}} \sum_{k=1}^{K} \sum_{i \in C_k} \| \mathbf{v}_i - \boldsymbol{\mu}_k \|^2
\]

where \(\mathbf{v}_i\) is grantee \(i\)'s velocity trajectory and
\(\boldsymbol{\mu}_k\) is cluster \(k\)'s centroid.

\subsubsection{Cluster Selection}\label{cluster-selection}

The number of clusters (K=3) is selected using:

\begin{enumerate}
\def\labelenumi{\arabic{enumi}.}
\tightlist
\item
  \textbf{Elbow method}: Plotting within-cluster sum of squares against
  K
\item
  \textbf{Silhouette scores}: Measuring cluster separation
\item
  \textbf{Interpretability}: Ensuring clusters have substantive meaning
\end{enumerate}

\subsubsection{Cluster Characterization}\label{cluster-characterization}

Clusters are characterized by:

\begin{itemize}
\tightlist
\item
  Mean velocity level (fast, moderate, slow)
\item
  Velocity variance (consistent vs.~variable)
\item
  Trajectory shape (accelerating, decelerating, flat)
\item
  Completion outcomes (median completion time, completion rate)
\end{itemize}

\subsection{B.5 Phase-Specific Velocity
Models}\label{b.5-phase-specific-velocity-models}

\subsubsection{Phase Definition}\label{phase-definition}

Programs are divided into terciles based on observed duration:

\[
\text{Phase}_t = \begin{cases}
\text{Early} & \text{if } t \leq T/3 \\
\text{Middle} & \text{if } T/3 < t \leq 2T/3 \\
\text{Late} & \text{if } t > 2T/3
\end{cases}
\]

\subsubsection{Multivariate Model}\label{multivariate-model}

Phase-specific effects are estimated by including all three phase
velocities simultaneously:

\[
\log h(t) = \log h_0(t) + \beta_E \text{Velocity}_{\text{Early}} + \beta_M \text{Velocity}_{\text{Mid}} + \beta_L \text{Velocity}_{\text{Late}}
\]

This specification identifies the independent contribution of each
phase, controlling for the others.

\subsubsection{Sensitivity Analysis}\label{sensitivity-analysis}

Alternative phase definitions are tested:

\begin{itemize}
\tightlist
\item
  \textbf{Quartiles}: Four equal phases
\item
  \textbf{Fixed intervals}: Years 1-3, 4-7, 8+
\item
  \textbf{Milestone-based}: Pre-50\%, 50-80\%, 80-95\%
\end{itemize}

Results are qualitatively similar across specifications.

\subsection{B.6 Meta-Analysis of Velocity
Estimates}\label{b.6-meta-analysis-of-velocity-estimates}

\subsubsection{Estimate Extraction}\label{estimate-extraction}

From each stratified model, we extract:

\begin{itemize}
\tightlist
\item
  Point estimate: \(\hat{\beta}\) (log hazard ratio)
\item
  Standard error: \(\text{SE}(\hat{\beta})\)
\item
  Sample size: \(n_k\)
\item
  Number of events: \(d_k\)
\end{itemize}

\subsubsection{Summary Statistics}\label{summary-statistics}

Rather than pooled estimation (which would require independence
assumptions), we compute:

\begin{itemize}
\tightlist
\item
  \textbf{Median HR}: Central tendency of effect estimates
\item
  \textbf{Range}: Minimum to maximum HR
\item
  \textbf{Interquartile range}: Middle 50\% of estimates
\item
  \textbf{Percent significant}: Proportion with p \textless{} 0.05
\item
  \textbf{Percent accelerating}: Proportion with HR \textgreater{} 1
\end{itemize}

\subsubsection{Forest Plot}\label{forest-plot}

Visual display shows each estimate with 95\% confidence intervals on a
log scale, organized by analysis type. Reference line at HR=1 indicates
null effect.

\subsection{B.7 Sensitivity Analyses}\label{b.7-sensitivity-analyses}

\subsubsection{Threshold Sensitivity}\label{threshold-sensitivity}

Completion thresholds other than 95\% are tested:

\begin{longtable}[]{@{}llll@{}}
\toprule\noalign{}
Threshold & N Events & Median HR & p-value Range \\
\midrule\noalign{}
\endhead
\bottomrule\noalign{}
\endlastfoot
90\% & 58 & 4.21 & 0.008-0.156 \\
95\% & 40 & 4.60 & 0.002-0.891 \\
99\% & 25 & 5.12 & 0.001-0.743 \\
\end{longtable}

Results are robust to threshold choice.

\subsubsection{QA Flag Exclusion}\label{qa-flag-exclusion}

Sensitivity analyses excluding observations flagged for data quality
issues:

\begin{longtable}[]{@{}lll@{}}
\toprule\noalign{}
Exclusion Criterion & N Excluded & HR Change \\
\midrule\noalign{}
\endhead
\bottomrule\noalign{}
\endlastfoot
Extreme velocity & 3 & +4.2\% \\
Large obligation adjustment & 12 & -6.1\% \\
Negative expenditure & 8 & +2.8\% \\
Any flag & 18 & -3.4\% \\
\end{longtable}

Results are substantively unchanged by exclusions.

\subsubsection{Winsorization Thresholds}\label{winsorization-thresholds}

Alternative winsorization cutoffs:

\begin{longtable}[]{@{}lll@{}}
\toprule\noalign{}
Cutoff & Median HR & 95\% CI \\
\midrule\noalign{}
\endhead
\bottomrule\noalign{}
\endlastfoot
1\%/99\% (baseline) & 4.60 & 1.42-14.89 \\
2.5\%/97.5\% & 4.43 & 1.38-14.21 \\
5\%/95\% & 4.28 & 1.31-13.98 \\
\end{longtable}

Effects are modestly attenuated with stricter winsorization but remain
statistically significant.

\subsection{B.8 Software and
Reproducibility}\label{b.8-software-and-reproducibility}

\subsubsection{Computational
Environment}\label{computational-environment}

\begin{itemize}
\tightlist
\item
  \textbf{Python 3.11} with lifelines package for survival analysis
\item
  \textbf{scikit-learn} for K-means clustering
\item
  \textbf{pandas/numpy} for data manipulation
\item
  \textbf{matplotlib/seaborn} for visualization
\end{itemize}

\subsubsection{Code Availability}\label{code-availability}

Analysis scripts are available in the project repository:

\begin{itemize}
\tightlist
\item
  \texttt{run\_phase\_specific\_analysis.py}: Phase velocity models
\item
  \texttt{run\_trajectory\_clustering.py}: K-means clustering
\item
  \texttt{run\_learning\_curves.py}: Experience stratification
\item
  \texttt{run\_program\_type\_analysis.py}: Program type heterogeneity
\item
  \texttt{run\_disaster\_context\_analysis.py}: Disaster context
  heterogeneity
\item
  \texttt{run\_meta\_analysis.py}: Summary statistics and forest plots
\end{itemize}

\subsubsection{Reproducibility}\label{reproducibility}

All analyses are reproducible from the standardized panel data:

\begin{Shaded}
\begin{Highlighting}[]
\ExtensionTok{python}\NormalTok{ src/pipeline.py ingest\_data}
\ExtensionTok{python}\NormalTok{ src/pipeline.py standardize\_data}
\ExtensionTok{python}\NormalTok{ src/pipeline.py build\_panel}
\ExtensionTok{python}\NormalTok{ src/pipeline.py build\_features\_std}
\ExtensionTok{python}\NormalTok{ run\_meta\_analysis.py}
\end{Highlighting}
\end{Shaded}





\end{document}
